\documentclass[11pt]{article}
\usepackage{manual_docs_style}
\externaldocument{build/cheap_filter_metrics_stage}[cheap_filter_metrics_stage.pdf]
\externaldocument{build/surrogate_training_stage}[surrogate_training_stage.pdf]
\externaldocument{build/candidate_generation}[candidate_generation.pdf]

\begin{document}

\docTitle{Stage: Surrogate Confidence Scoring}
\phantomsection\label{docstart:surrogate_confidence_scoring_stage}

This stage applies a frozen surrogate to the production candidate pool after cheap filtering.
It writes per-run confidence records used by dataset selection.

\section*{Main Driver}

\begin{DocCode}
env/bin/python workflows/surrogate/score_surrogate.py \
  --model SIMULATOR \
  --policy PRODUCTION_POLICY_ID \
  --surrogate-id SURROGATE_ID \
  [--noise-level none|low|high|all]
\end{DocCode}

\section*{Inputs}

\begin{itemize}
\item \code{plans/<production_policy_id>/run_nodes.jsonl}
\item \code{plans/<production_policy_id>/run_edges.jsonl}
\item \code{runs/model_record.json}
\item \code{runs/<run_id>/data.parquet}
\item \code{metrics/<production_policy_id>/cheap_filter_metrics.jsonl}
\item \code{surrogates/<surrogate_id>/model.pt}
\item \code{surrogates/<surrogate_id>/thresholds.json}
\item \code{experiment_config.json}
\item \code{noise_adder.py}
\end{itemize}

\section*{Output}

The stage writes:

\begin{DocCode}
models/simulink/<simulator_id>/metrics/<production_policy_id>/surrogate_scores/<surrogate_id>.jsonl
\end{DocCode}

Records are per \code{run_id} and per noise level because a sample that is recognizable with low noise may not be recognizable with high noise.

\begin{DocCode}
{
  "record_type": "surrogate_score",
  "policy_id": "production_v1",
  "surrogate_id": "surrogate_ball_drop_v1",
  "run_id": "run_7fd9daa0055fb3f412ef5e0ab8d4c5ba",
  "family_id": "fam_de4d7c0a5335e9005510cc2442d0971e",
  "kind": "intervention",
  "ground_truth_label": "gravity",
  "confidence_correct_label": [(0,0.99,...., 10,0.4)]
  "confidence_argmax_label": [(0,0.99,...., 10,0.4)]
  "argmax_label":[(0,"mass",..."10"",gravity")]
}
\end{DocCode}

In these tuple-valued fields, the first element is the noise multiplier relative to the simulator's \code{low} noise profile.
A multiplier of \code{0} means no added observation noise, \code{1} means the base \code{low} noise magnitude, and \code{10} means ten times the \code{low} noise magnitude.
\code{confidence_correct_label} records the surrogate confidence assigned to \code{ground_truth_label} at each multiplier.
\code{confidence_argmax_label} records the confidence of the surrogate's argmax label at each multiplier, and \code{argmax_label} records that predicted label.

\section*{Pass Rule}

\code{surrogate_filter_pass} is true only if:

\begin{enumerate}
\item \code{cheap_filter_pass == true};
\item \code{predicted_label == ground_truth_label};
\item \code{true_label_confidence} is at least the applicable global or per-class threshold;
\item \code{confidence_margin} is at least the configured margin threshold, when a margin threshold is enabled.
\end{enumerate}

Using \code{max_confidence} alone is not sufficient.
A confidently wrong sample must be rejected.
Selection must use confidence in the simulator ground-truth label.

\end{document}
